\documentclass[11pt]{article}
\usepackage[margin=1in]{geometry}
\usepackage{amsmath, amssymb, amsfonts}
\usepackage{graphicx}
\usepackage{booktabs}
\usepackage{hyperref}
\usepackage{enumitem}
\usepackage{float}
\usepackage{subcaption}
\usepackage{xcolor}
\usepackage{listings}
\usepackage{fancyhdr}

\lstset{
  basicstyle=\ttfamily\small,
  breaklines=true,
  frame=single,
  language=Python,
  keywordstyle=\color{blue},
  commentstyle=\color{gray},
  stringstyle=\color{red},
}

\pagestyle{fancy}
\fancyhf{}
\rhead{CS336 Assignment 2}
\lhead{Systems and Parallelism}
\cfoot{\thepage}

\title{\textbf{CS336 Assignment 2: Systems and Parallelism} \\ Writeup}
\author{Your Name \\ SUNet ID}
\date{Spring 2026}

\begin{document}
\maketitle
\tableofcontents
\newpage

%% ============================================================
\section{Profiling and Benchmarking}
%% ============================================================

% ----------------------------------------------------------
\subsection{End-to-End Benchmarking}
\textbf{Problem (benchmarking\_script)}

\paragraph{(b) Timings for forward, backward, and optimizer step.}

% TODO: Insert your timing table here
% \begin{table}[H]
% \centering
% \caption{End-to-end benchmarking results (5 warmup, 10 measurement steps).}
% \begin{tabular}{lrrr}
% \toprule
% Model Size & Forward (ms) & Backward (ms) & Optimizer Step (ms) \\
% \midrule
% small  & & & \\
% medium & & & \\
% large  & & & \\
% xl     & & & \\
% \bottomrule
% \end{tabular}
% \end{table}

\paragraph{(c) Effect of removing warm-up steps.}


% ----------------------------------------------------------
\subsection{Nsight Systems Profiling}
\textbf{Problem (nsys\_profile)}

\paragraph{(a) Total time on forward pass.}


\paragraph{(b) Most time-consuming CUDA kernel.}


\paragraph{(c) Non-matmul kernels with non-trivial runtime.}


\paragraph{(d) Full training step profile comparison.}


\paragraph{(e) Softmax vs.\ matrix multiplication runtime comparison.}


% ----------------------------------------------------------
\subsection{Mixed Precision}

\subsubsection{Mixed-Precision Accumulation}
\textbf{Problem (mixed\_precision\_accumulation)}


\subsubsection{Benchmarking Mixed Precision}
\textbf{Problem (benchmarking\_mixed\_precision)}

\paragraph{(a) Data types under FP16 autocasting.}

% TODO: List data types for each component

\paragraph{(b) Layer normalization and BF16.}


\paragraph{(c) Mixed precision benchmarking results.}

% TODO: Insert timing table with/without mixed precision

% ----------------------------------------------------------
\subsection{Memory Profiling}
\textbf{Problem (memory\_profiling)}

\paragraph{(a) Memory timelines.}

% TODO: Insert two images from memory_viz tool

\paragraph{(b) Peak memory usage.}

% TODO: Insert table with peak memory per context length

\paragraph{(c) Mixed precision effect on memory.}


\paragraph{(d) Size of activation tensor in residual stream.}


\paragraph{(e) Largest allocations in memory timeline.}


\paragraph{(f) Memory saved for backward per TransformerBlock.}

% TODO: Insert Nsight screenshots


%% ============================================================
\section{Single-GPU Memory}
%% ============================================================

\textbf{Problem (gradient\_checkpointing)}

\paragraph{(a) Optimal checkpointing strategy.}

% TODO: Describe strategy, asymptotic memory, code sketch

\paragraph{(b) Best single-level checkpointing for xl model.}

% TODO: Describe reasoning and measured peak memory


%% ============================================================
\section{GPU Kernels}
%% ============================================================

% ----------------------------------------------------------
\subsection{PyTorch Attention Benchmarking}
\textbf{Problem (pytorch\_attention)}

\paragraph{(a) Attention benchmarking at different scales.}

% TODO: Insert timing table and memory analysis

% ----------------------------------------------------------
\subsection{Torch Compile}
\textbf{Problem (torch\_compile)}

\paragraph{(a) Compiled vs.\ uncompiled attention.}

% TODO: Insert comparison table

\paragraph{(b) Compiled full Transformer model.}

% TODO: Insert comparison table

% ----------------------------------------------------------
\subsection{FlashAttention-2 Benchmarking}
\textbf{Problem (flash\_benchmarking)}

\paragraph{(a) Triton FlashAttention-2 vs.\ PyTorch attention.}

% TODO: Insert comparison table


%% ============================================================
\section{Distributed Data Parallel Training}
%% ============================================================

% ----------------------------------------------------------
\subsection{Distributed Communication}
\textbf{Problem (distributed\_communication\_single\_node)}

% TODO: Insert plots/tables and commentary

% ----------------------------------------------------------
\subsection{Na\"ive DDP Benchmarking}
\textbf{Problem (naive\_ddp\_benchmarking)}

% TODO: Description and measured timings

% ----------------------------------------------------------
\subsection{Minimal DDP with Flat Gradients}
\textbf{Problem (minimal\_ddp\_flat\_benchmarking)}

% TODO: Measured timings and comparison

% ----------------------------------------------------------
\subsection{DDP with Overlapping Individual Parameters}
\textbf{Problem (ddp\_overlap\_individual\_parameters\_benchmarking)}

\paragraph{(a) Benchmarking overlapped DDP.}

% TODO: Measured timings and comparison

\paragraph{(b) Nsight profiler comparison.}

% TODO: Insert 2 Nsight screenshots


%% ============================================================
\section{Optimizer State Sharding}
%% ============================================================

\textbf{Problem (optimizer\_state\_sharding\_accounting)}

\paragraph{(a) Peak memory usage with and without sharding.}

% TODO: Memory results and breakdown

\paragraph{(b) Training speed with and without sharding.}

% TODO: Timing results

\paragraph{(c) Comparison with ZeRO Stage 1.}



%% ============================================================
\section{Fully-Sharded Data Parallel}
%% ============================================================

\textbf{Problem (fsdp\_accounting)}

\paragraph{(a) Expected memory savings from FSDP.}


\paragraph{(b) All-gather timing analysis.}

% TODO: Insert Nsight screenshots and timings


%% ============================================================
\section{Analyzing Parallelism Strategies}
%% ============================================================

% ----------------------------------------------------------
\subsection{Communication Primitives}
\textbf{Problem (alternate\_ring\_all\_reduce)}

% TODO: Answer in terms of S, N, W with justification

% ----------------------------------------------------------
\subsection{Data Parallel Calculations}
\textbf{Problem (data\_parallel\_calcs)}

\paragraph{(a) FLOPs for backward pass with $N_{\text{DP}}$ data parallelism.}


\paragraph{(b) Communication time for backward pass.}


\paragraph{(c) Maximum $N_{\text{DP}}$ before communication bottleneck.}


% ----------------------------------------------------------
\subsection{FSDP Calculations}
\textbf{Problem (fsdp\_calcs)}

\paragraph{(a) FLOPs for forward and backward pass with FSDP.}


\paragraph{(b) Communication time for forward and backward pass.}


\paragraph{(c) Maximum $N_{\text{FSDP}}$ before communication bottleneck.}


% ----------------------------------------------------------
\subsection{Tensor Parallel Calculations}
\textbf{Problem (tp\_calcs)}

\paragraph{(a) Backward pass equations for TP.}

% TODO: Series of equations

\paragraph{(b) FLOPs for forward and backward pass with TP.}


\paragraph{(c) Communication time for forward and backward pass.}


\paragraph{(d) Maximum $N_{\text{TP}}$ before communication bottleneck.}


% ----------------------------------------------------------
\subsection{2D Parallelism (FSDP + TP)}
\textbf{Problem (fsdp\_tp\_calcs)}

\paragraph{(a) FLOPs for forward pass with 2D parallelism.}


\paragraph{(b) Communication time for forward pass.}


\paragraph{(c) Maximum $N$ with overlapped communication.}


\paragraph{(d) Maximum $N$ without overlapped communication.}



%% ============================================================
\section{Leaderboard}
%% ============================================================

\textbf{Problem (leaderboard)}

% TODO: Report best wall-clock time

\end{document}
